% Shared preamble of the proof overview documents  <Folder>/<Name>_overview.tex.
% Each overview starts with
%     % Shared preamble of the proof overview documents  <Folder>/<Name>_overview.tex.
% Each overview starts with
%     % Shared preamble of the proof overview documents  <Folder>/<Name>_overview.tex.
% Each overview starts with
%     \input{../../scripts/overview_preamble.tex}
% (for ThinMatching/: \input{../scripts/overview_preamble.tex}) and is compiled from its own folder by
% scripts/build_overviews.py.
\documentclass[10pt,a4paper]{article}
\usepackage[T1]{fontenc}
\usepackage[utf8]{inputenc}
\usepackage[margin=2.2cm]{geometry}
\usepackage{amsmath,amssymb,amsthm}
\usepackage{array,booktabs,longtable}
\usepackage[dvipsnames]{xcolor}
\usepackage{enumitem}
\usepackage{tikz}
\usetikzlibrary{arrows.meta,positioning,shapes.multipart,fit,calc}
\usepackage[hidelinks]{hyperref}
\setlist{itemsep=2pt,topsep=4pt}
\renewcommand{\arraystretch}{1.2}

% Lean names, files and paper references
\newcommand{\lean}[1]{\textcolor{blue!50!black}{\texttt{\detokenize{#1}}}}
\newcommand{\file}[1]{\textcolor{OliveGreen!70!black}{\texttt{\detokenize{#1}}}}
\newcommand{\paperref}[1]{\textsf{#1}}
\newcommand{\hyp}[1]{\textcolor{BrickRed}{\textbf{#1}}}

% Boxes for the dependency diagrams
\tikzset{
  res/.style={draw, rounded corners, fill=blue!6, align=center, inner sep=4pt, font=\small},
  main/.style={res, fill=blue!18, very thick},
  cited/.style={res, fill=red!10, dashed},
  shared/.style={res, fill=gray!12},
  dep/.style={-{Stealth[length=2.2mm]}, thick},
}

\newcommand{\overviewtitle}[2]{%
  \begin{center}{\Large\bfseries #1}\\[4pt]{\small Proof overview of the Lean~4/Mathlib formalisation}\\[2pt]
  {\small #2}\end{center}}

% Standard paragraph describing the paper
\newcommand{\paperversion}{\noindent\textbf{Paper.} J.~Breitung, T.~Gervens, M.~Grohe.
  \emph{Graph Similarity Meets Optimal Transport}.\\ File \file{Graph_Similarity_Meets_Optimal_Transport.pdf} in the
  repository; statement, equation, page and line numbers refer to this PDF. \file{docs/NUMBERING.md} lists the Lean
  files and declarations of all numbered statements.\par\smallskip}

% (for ThinMatching/: \input{../scripts/overview_preamble.tex}) and is compiled from its own folder by
% scripts/build_overviews.py.
\documentclass[10pt,a4paper]{article}
\usepackage[T1]{fontenc}
\usepackage[utf8]{inputenc}
\usepackage[margin=2.2cm]{geometry}
\usepackage{amsmath,amssymb,amsthm}
\usepackage{array,booktabs,longtable}
\usepackage[dvipsnames]{xcolor}
\usepackage{enumitem}
\usepackage{tikz}
\usetikzlibrary{arrows.meta,positioning,shapes.multipart,fit,calc}
\usepackage[hidelinks]{hyperref}
\setlist{itemsep=2pt,topsep=4pt}
\renewcommand{\arraystretch}{1.2}

% Lean names, files and paper references
\newcommand{\lean}[1]{\textcolor{blue!50!black}{\texttt{\detokenize{#1}}}}
\newcommand{\file}[1]{\textcolor{OliveGreen!70!black}{\texttt{\detokenize{#1}}}}
\newcommand{\paperref}[1]{\textsf{#1}}
\newcommand{\hyp}[1]{\textcolor{BrickRed}{\textbf{#1}}}

% Boxes for the dependency diagrams
\tikzset{
  res/.style={draw, rounded corners, fill=blue!6, align=center, inner sep=4pt, font=\small},
  main/.style={res, fill=blue!18, very thick},
  cited/.style={res, fill=red!10, dashed},
  shared/.style={res, fill=gray!12},
  dep/.style={-{Stealth[length=2.2mm]}, thick},
}

\newcommand{\overviewtitle}[2]{%
  \begin{center}{\Large\bfseries #1}\\[4pt]{\small Proof overview of the Lean~4/Mathlib formalisation}\\[2pt]
  {\small #2}\end{center}}

% Standard paragraph describing the paper
\newcommand{\paperversion}{\noindent\textbf{Paper.} J.~Breitung, T.~Gervens, M.~Grohe.
  \emph{Graph Similarity Meets Optimal Transport}.\\ File \file{Graph_Similarity_Meets_Optimal_Transport.pdf} in the
  repository; statement, equation, page and line numbers refer to this PDF. \file{docs/NUMBERING.md} lists the Lean
  files and declarations of all numbered statements.\par\smallskip}

% (for ThinMatching/: \input{../scripts/overview_preamble.tex}) and is compiled from its own folder by
% scripts/build_overviews.py.
\documentclass[10pt,a4paper]{article}
\usepackage[T1]{fontenc}
\usepackage[utf8]{inputenc}
\usepackage[margin=2.2cm]{geometry}
\usepackage{amsmath,amssymb,amsthm}
\usepackage{array,booktabs,longtable}
\usepackage[dvipsnames]{xcolor}
\usepackage{enumitem}
\usepackage{tikz}
\usetikzlibrary{arrows.meta,positioning,shapes.multipart,fit,calc}
\usepackage[hidelinks]{hyperref}
\setlist{itemsep=2pt,topsep=4pt}
\renewcommand{\arraystretch}{1.2}

% Lean names, files and paper references
\newcommand{\lean}[1]{\textcolor{blue!50!black}{\texttt{\detokenize{#1}}}}
\newcommand{\file}[1]{\textcolor{OliveGreen!70!black}{\texttt{\detokenize{#1}}}}
\newcommand{\paperref}[1]{\textsf{#1}}
\newcommand{\hyp}[1]{\textcolor{BrickRed}{\textbf{#1}}}

% Boxes for the dependency diagrams
\tikzset{
  res/.style={draw, rounded corners, fill=blue!6, align=center, inner sep=4pt, font=\small},
  main/.style={res, fill=blue!18, very thick},
  cited/.style={res, fill=red!10, dashed},
  shared/.style={res, fill=gray!12},
  dep/.style={-{Stealth[length=2.2mm]}, thick},
}

\newcommand{\overviewtitle}[2]{%
  \begin{center}{\Large\bfseries #1}\\[4pt]{\small Proof overview of the Lean~4/Mathlib formalisation}\\[2pt]
  {\small #2}\end{center}}

% Standard paragraph describing the paper
\newcommand{\paperversion}{\noindent\textbf{Paper.} J.~Breitung, T.~Gervens, M.~Grohe.
  \emph{Graph Similarity Meets Optimal Transport}.\\ File \file{Graph_Similarity_Meets_Optimal_Transport.pdf} in the
  repository; statement, equation, page and line numbers refer to this PDF. \file{docs/NUMBERING.md} lists the Lean
  files and declarations of all numbered statements.\par\smallskip}

\begin{document}
\overviewtitle{Theorem 17: inapproximability of the optimal-transport operator $p$-distances}{Folder \file{GsimOT/Theorem18_OperatorNormHardness}}
\paperversion

\section{Statement}

\paragraph{Paper.} Theorem~17 (Section~4, p.~12, l.~432--433; proof in App.~G, pp.~54--57,
l.~1710--1804): \emph{for $1\le p\le\infty$ and any $\varepsilon>0$, unless $\mathrm{P}=\mathrm{NP}$, there
is no polynomial-time approximation algorithm computing $d^{\mathrm{OT}}_p$ with multiplicative error
guarantee $2^{b(p)}-\varepsilon$}, where $b(p)=\max\{1/p,1-1/p\}$ and $b(\infty)=1$ (l.~431). It gives
Theorem~4\,(2) (p.~4, l.~135--137); the restriction to maximum degree two for $p=1$ is Theorem~16. The tools
are the aggregated matrix $A^Q_{G,H}=A_G-QA_HQ^\top$ of Eq.~(45) (p.~53),
Lemma~53 (p.~53, $\|A^Q_{G,H}\|_p\le\|B_p(G,H,Q)\|_p$) and Corollary~54 (p.~54, a column with $k$ entries
of size $\ge\varepsilon$ gives $\|A^Q_{G,H}\|_p\ge k^{b(p)}\varepsilon$). The reduction from Three-Partition
takes the paths and cycles of Theorem~16 and attaches $M$ leaves to every vertex; YES-instances give
$d^{\mathrm{OT}}_p(G,H)\le1$, NO-instances $d^{\mathrm{OT}}_p(G,H)>2^{b(p)}-\varepsilon^*$ (p.~54, l.~1728--1730).

\paragraph{Lean.} Only the gap (the correctness of the reduction) is formalised. For an instance $I$, the
graphs are $G=$ \lean{I.graphG M} $=$ \lean{addLeaves (pathUnion I.a) (fun _ => M)} and
$H=$ \lean{I.graphH M} $=$ \lean{addLeaves (cycleUnion I.m I.A) (fun _ => M)}, and $1\le p\le\infty$
(\lean{p : ENNReal} with \lean{Fact (1 <= p)}).
\begin{itemize}
\item \lean{GsimOT.Hardness.ThreePartition.dOT_le_one_of_isYes}: if $a_i\ge2$ for all $i$ and $I$ is a
  YES-instance, then $d^{\mathrm{OT}}_p(G,H)\le1$ (for every $M$);
\item \lean{GsimOT.Hardness.ThreePartition.lt_dOT_of_not_isYes}: if $a_i\ge3$ for all $i$, $I$ is a
  NO-instance, $0<\varepsilon\le1$ and $M>\max\{2(4/\eta)^2+1,\ 2(2/\varepsilon')^2+1\}$ with
  $\eta=\varepsilon/8$, $\varepsilon'=(\varepsilon/8)/(mA-1)$ (Eq.~(46)), then
  $2^{b(p)}-\varepsilon<d^{\mathrm{OT}}_p(G,H)$;
\item \lean{GsimOT.Hardness.ThreePartition.dOT_shift_gap}: both, for every instance, after the
  w.l.o.g.\ shift (add $3$ to each $a_i$ and $9$ to $A$; the paper assumes $a_i\ge3$ w.l.o.g., l.~1712).
\end{itemize}
$d^{\mathrm{OT}}_p=$ \lean{dOT p p} is the infimum over $\mathrm{DS}(V,W)$ of $\|B_p(G,H,Q)\|_p$
(\lean{Bmat p p}, \lean{opNorm p p}; \file{Common/OTDistances.lean}, \file{Common/OperatorNorm.lean}).
The Lean statement is the paper's gap, for all $0<\varepsilon\le1$; this range contains the paper's
$0<\varepsilon^*<\sqrt2-1$ (l.~1728).

\section{Where the proof lives}
\begin{longtable}{>{\raggedright\arraybackslash}p{5.2cm} >{\raggedright\arraybackslash}p{10.6cm}}
\toprule \textbf{File} & \textbf{Content} \\ \midrule\endhead \bottomrule\endfoot
\file{OperatorPNorms.lean} & $b(p)$ (\lean{bExp}), Jensen (\lean{rpow_sum_mul_le}), Schur test
  (\lean{opNorm_le_one_of_sum_abs_le_one}), the column bound of Corollary~54 for symmetric matrices
  (\lean{rpow_bExp_mul_le_opNorm}). \\
\file{Lemma19_AggregatedMatrix.lean} & Eq.~(45) (\lean{aggMatrix}, \lean{aggMatrix_isSymm}) and Lemma~53
  (\lean{opNorm_aggMatrix_le}). \\
\file{Corollary20_ColumnBound.lean} & Corollary~54 (\lean{rpow_bExp_mul_le_opNorm_Bpq}, \lean{lt_opNorm_Bpq}). \\
\file{Claim54_LeafMass.lean} & the observation on $A^Q$ for (non-)edges (\lean{aggMatrix_of_adj},
  \lean{aggMatrix_of_not_adj}), $|N_G(v)|\ge M+1$, Claim~55 (\lean{claim54}). \\
\file{Claim55_SplitMass.lean} & Claim~56 (\lean{claim55}). \\
\file{Theorem18.lean} & the graphs $G,H$, the YES-direction, the final argument (NO-direction), the shifted version. \\
\file{Common/ThreePartition.lean}, \file{Common/PathsAndCycles.lean}, \file{Common/Leaves.lean} &
  instances, layout bijection, counting step \lean{isYes_of_isDS}, paths, cycles, leaves. \\
\end{longtable}

\section{Proof layout}
\begin{center}
\begin{tikzpicture}
  \node[main] (gap) at (0,0) {Theorem 17 (any instance)\\\lean{dOT_shift_gap}};
  \node[main] (yes) at (-4.4,-1.9) {YES: $d^{\mathrm{OT}}_p\le1$\\\lean{dOT_le_one_of_isYes}};
  \node[main] (no) at (3.4,-1.9) {NO: $d^{\mathrm{OT}}_p>2^{b(p)}-\varepsilon$\\\lean{lt_dOT_of_not_isYes}};
  \node[res] (schur) at (-4.6,-4.1) {Schur test\\\lean{opNorm_Bpq_permMat_le_one}};
  \node[res] (c55) at (-1.3,-4.1) {Claim 55\\\lean{claim54}};
  \node[res] (c56) at (1.4,-4.1) {Claim 56\\\lean{claim55}};
  \node[shared] (count) at (5.6,-4.1) {counting step\\\lean{isYes_of_isDS}};
  \node[res] (cor) at (1.0,-6.2) {Corollary 54\\\lean{rpow_bExp_mul_le_opNorm_Bpq}\\\lean{lt_opNorm_Bpq}};
  \node[res] (lem) at (-2.6,-8.3) {Lemma 53\\\lean{opNorm_aggMatrix_le}};
  \node[res] (col) at (3.6,-8.3) {column bound, symmetric $A$\\\lean{rpow_bExp_mul_le_opNorm}};
  \node[shared] (lay) at (-4.6,-6.2) {layout bijection\\\lean{layoutEquiv}};
  \node[shared] (norm) at (-2.6,-10.2) {$\ell^p$ / operator norms (Common)\\\lean{opNorm}, \lean{Bmat}};
  \draw[dep] (gap) -- (yes); \draw[dep] (gap) -- (no);
  \draw[dep] (yes) -- (schur); \draw[dep] (yes.west) to[out=200,in=170,looseness=1.4] (lay.west);
  \draw[dep] (no) -- (c55); \draw[dep] (no) -- (c56); \draw[dep] (no) -- (count); \draw[dep] (no) -- (cor.north east);
  \draw[dep] (c55) -- (cor); \draw[dep] (c56) -- (cor);
  \draw[dep] (cor) -- (lem); \draw[dep] (cor) -- (col);
  \draw[dep] (lem) -- (norm);
\end{tikzpicture}
\end{center}
Blue: this folder; grey: \file{GsimOT/Common}. Nothing is cited.

\section{The proof step by step}

\paragraph{Eq.~(45) and Lemma~53 (pp.~52--54, l.~1671--1697).} $A^Q_{G,H}(v,v')=\sum_{w,w'}Q(v,w)D(v,w,v',w')Q(v',w')$
is \lean{aggMatrix}; it is symmetric (\lean{aggMatrix_isSymm}, used in Corollary~54). Lemma~53 is
\lean{opNorm_aggMatrix_le}. The paper picks optimal test vectors $x,y$ with Lemma~19 (duality, App.~A, p.~14) and lifts them to
$\tilde x(v,w)=Q(v,w)^{1/p}x(v)$, $\tilde y(v,w)=Q(v,w)^{1-1/p}y(v)$, so that $y^\top A^Qx=\tilde y^\top B_p\tilde x$.
The formalisation uses the same $\tilde x$ but no $\tilde y$: it writes
$(A^Qx)(v)=\sum_wQ(v,w)^{1-1/p}(B_p\tilde x)(v,w)$ and bounds $\|A^Qx\|_p\le\|B_p\tilde x\|_p$ by Jensen
(\lean{rpow_sum_mul_le}), since the rows of $Q$ sum to $1$. The convention $0^0=0$ is \lean{powZ}.

\paragraph{Corollary~54 (p.~54, l.~1699--1709).} If column $v$ of $A^Q$ has $k$ entries of absolute value
$\ge\varepsilon$, then $\|A^Q\|_p\ge k^{b(p)}\varepsilon$ (\lean{rpow_bExp_mul_le_opNorm}). The bound
$k^{1/p}\varepsilon$ comes from $e_v$, as in the paper; for $k^{1-1/p}\varepsilon$ the paper uses
$\|A\|_p=\|A^\top\|_{p^*}$ and symmetry, the formalisation tests $A$ directly with a normalised sign vector on
the $k$ rows, which is the same computation. With Lemma~53 this gives \lean{rpow_bExp_mul_le_opNorm_Bpq}, and
\lean{lt_opNorm_Bpq} is the ``in particular'' ($k>(c/\varepsilon)^{1/b(p)}\Rightarrow\|B_p\|_p>c$).

\paragraph{Construction (p.~54, l.~1710--1727).} $G$, $H$ are \lean{graphG}, \lean{graphH}
(\lean{addLeaves}); base vertices are \lean{Sum.inl}, leaves \lean{Sum.inr}. The w.l.o.g.\ $a_i\ge3$ is
\lean{ThreePartition.shift}; $A\ge7$ is derived in \lean{lt_dOT_of_not_isYes} (\lean{hA}). The observation
$A^Q(x,x')=\sum_{yy'\notin E(H)}Q(x,y)Q(x',y')$ for $xx'\in E(G)$ is \lean{aggMatrix_of_adj}; its
counterpart for non-edges is \lean{aggMatrix_of_not_adj}.

\paragraph{YES-instances (p.~55, l.~1733--1737).} The layout bijection $\tau$ of three paths per cycle
(\lean{layoutEquiv}, \lean{layout_adj_cycleUnion}) is extended to the leaves by
\lean{Equiv.sumCongr} and \lean{Equiv.sigmaCongrLeft}. The paper observes that $A_H-A^\pi_G$ is the
adjacency matrix of $3m$ disjoint edges, so its norm is $1$. The formalisation shows that every vertex has
at most one mismatch and bounds $\|B_p(G,H,P_\pi)\|_p\le1$ by the Schur test
(\lean{opNorm_Bpq_permMat_le_one}, \lean{opNorm_le_one_of_sum_abs_le_one}).

\paragraph{Claim~55 (p.~55, l.~1741--1759).} \lean{claim54}: if $Q(v,L_H)\ge\eta$ and
$M>2(4/\eta)^2+1$, then $\|B_p\|_p>2^{b(p)}$. The case split ``for all $w$, $Q(v,L(w))\le\eta/2$''
(\lean{hcase}) and the two lower bounds $\eta/2$ (Case~1) and $\frac\eta2(1-\delta)$ with $\delta=1/2$ (Case~2)
follow the paper line by line; the at most $\lfloor1/\delta\rfloor=2$ heavy vertices are counted by
\lean{card_mul_le_of_lt_mass}. The paper's size condition $M-1>2(4/\eta)^2$ is matched to Corollary~54 by
\lean{rpow_inv_bExp_le}: $(2^{b(p)}a)^{1/b(p)}\le2a^2$ for $a\ge1$.

\paragraph{Claim~56 (pp.~55--56, l.~1760--1773).} \lean{claim55}: if $Q(v,w),Q(v,w')\ge\varepsilon$ for distinct
base vertices with at most one common neighbour and $M>2(2/\varepsilon)^2+1$, then $\|B_p\|_p>2^{b(p)}$.
The paper's ``$|I|\le1$ since $|V(C_j)|=A\ge5$'' (l.~1762) is the hypothesis \lean{hI}, discharged by
\lean{card_common_adj_cycleUnion_le_one}.

\paragraph{Final argument (pp.~56--57, l.~1774--1804).} In \lean{lt_dOT_of_not_isYes}: if Claim~55 or
Claim~56 applies, done (\lean{hA'}, \lean{hB}). Otherwise every base vertex $v$ sends mass
$\ge1-\varepsilon/8-\varepsilon/8=1-\varepsilon/4$ to a single base vertex $f(v)$ (\lean{hf}, the paper's
$Q(v,w)\ge1-\varepsilon-\eta$); since $1-\varepsilon/4>1/2$, $f$ is injective, hence the bijection $\tau$
(\lean{hf_inj}, \lean{Equiv.ofBijective}). As $I$ is a NO-instance, some edge $vv'$ of a path is mapped into
two different cycles (\lean{hcross}, via the counting step \lean{isYes_of_isDS}); with $a_i\ge3$ the vertex
$v$ has a second neighbour $v''$ (\lean{exists_other_adj}). Either $\tau(v'')\not\sim\tau(v)$ (an entry
$\ge(1-\varepsilon/4)^2$ by \lean{hentry_adj}) or $\tau(v)$ has a cycle neighbour $w^*\ne\tau(v'')$ whose
preimage is not adjacent to $v$ (an entry $\le-(1-\varepsilon/4)^2$, via \lean{aggMatrix_of_not_adj}). Two
entries in column $v$ and Corollary~54 with $k=2$ give $\|B_p\|_p\ge2^{b(p)}(1-\varepsilon/4)^2>2^{b(p)}-\varepsilon$
(\lean{hL}, using $2^{b(p)}\le2$). The infimum over $Q$ is handled by \lean{le_iInf_isDS}.

\section{Assumptions and trusted base}
No cited results and no hypotheses beyond those of the statement. Only the standard axioms \lean{propext},
\lean{Classical.choice}, \lean{Quot.sound}. The complexity-theoretic part (polynomial size of the reduction,
NP-hardness of Three-Partition) is not formalised.

\section{Notes on the formalisation}
\begin{itemize}
\item Lemma~53 is proved by Jensen instead of the duality of Lemma~19; Corollary~54 uses a sign vector
  instead of $\|A\|_p=\|A^\top\|_{p^*}$ (see above). The bounds are the same.
\item YES-direction: the Schur test (bound $\le1$) instead of the computation of the norm ($=1$);
  \lean{dOT_le_one_of_isYes} assumes $a_i\ge2$.
\item $\varepsilon$ ranges over $(0,1]$, which contains the paper's range $(0,\sqrt2-1)$ for $\varepsilon^*$
  (l.~1728); the paper uses $\varepsilon^*<\sqrt2-1$ to conclude $2^{b(p)}-\varepsilon^*>1$ (l.~1729--1730).
\item Remark~62 (other operator $p\to q$-norms, p.~62, l.~1989--1998) is not formalised.
\end{itemize}

\section{How to check this result}
\begin{itemize}
\item \texttt{lake build GsimOT.Theorem18\_OperatorNormHardness.Theorem18}
\item \texttt{lake env lean scripts/CheckAxioms.lean} prints the axioms of
  \lean{GsimOT.Hardness.ThreePartition.dOT_le_one_of_isYes} and \lean{GsimOT.Hardness.ThreePartition.lt_dOT_of_not_isYes}.
\item The table \file{Theorem18_correspondence.pdf} in this folder lists every paper item with its Lean
  declaration and line.
\end{itemize}
\end{document}
